\begin{table}[t]
\centering\small\setlength{\tabcolsep}{4.5pt}
\begin{tabular}{lrrrr}
\toprule
Method & Return & Plan s & Unknown & Max debt\\
\midrule
Depth two & 4.197 & 8.344 & 30.7\% & 0\\
Free info. & 3.803 & 0.940 & 0.0\% & 0\\
Bounded opt. & 4.392 & 1.560 & 0.1\% & 9\\
V4 frozen & 4.391 & 0.010 & --- & 6\\
Full history & 4.432 & 0.036 & --- & 6\\
\bottomrule
\end{tabular}
\caption{Frozen v6 holdout: mean episode return, total planning seconds per 24-episode stream, and maximum completed-prefix anchor deficit. Only the first three arms share line metering and the one-second ceiling. Full-history optimism has no risk ledger. All ranker weights stay frozen. Timing includes measurement overhead.}
\label{tab:v6}
\end{table}
The primary mean difference is $+0.39375$ reward per episode, with paired 95\% Student interval $[+0.15845,+0.62905]$ over 40 independent streams. The depth-two selector uses 8.88$\times$ the cheap control's measured planning time. The holdout contains 40 distinct tables; 0 streams match a development table (0 distinct overlapping tables). This is same-family validation, not new-domain generalization.
\begin{figure}[t]
\centering
\includegraphics[width=\columnwidth]{figures/v6_reward_cost.pdf}
\caption{Reward and measured planning cost for the three commonly metered v6 arms. Reward intervals are marginal 95\% Student intervals across streams; the primary analysis uses paired differences. Common ceilings do not imply equal realized computation.}
\label{fig:v6cost}
\end{figure}
